\textbf{1. Kiến thức cần nhớ}
\begin{itemize}
\item \textbf{Elip ($E$):} $\frac{x^2}{a^2} + \frac{y^2}{b^2} = 1$ ($a > b > 0$).
  \begin{itemize}
  \item $c = \sqrt{a^2 - b^2}$.
  \item Tiêu điểm: $F_1(-c; 0), F_2(c; 0)$[cite: 1]. Tiêu cự: $2c$.
  \item Độ dài trục lớn: $2a$. Độ dài trục nhỏ: $2b$.
  \end{itemize}
\item \textbf{Hypebol ($H$):} $\frac{x^2}{a^2} - \frac{y^2}{b^2} = 1$ ($a, b > 0$)[cite: 1].
  \begin{itemize}
  \item $c = \sqrt{a^2 + b^2}$[cite: 1].
  \item Tiêu điểm: $F_1(-c; 0), F_2(c; 0)$[cite: 1]. Tiêu cự: $2c$[cite: 1].
  \item Độ dài trục thực: $2a$. Độ dài trục ảo: $2b$.
  \end{itemize}
\item \textbf{Parabol ($P$):} $y^2 = 2px$ ($p > 0$)[cite: 1].
  \begin{itemize}
  \item Tiêu điểm: $F\left(\frac{p}{2}; 0\right)$[cite: 1]. Tham số tiêu: $p$.
  \item Phương trình đường chuẩn $\Delta: x = -\frac{p}{2}$[cite: 1].
  \end{itemize}
\end{itemize}

\textbf{2. Phương pháp giải các dạng}

%% DANG 0H9?5-1
\textbf{Dạng 1. Xác định các yếu tố của elip}
\begin{enumerate}
\item Từ phương trình $\frac{x^2}{a^2} + \frac{y^2}{b^2} = 1$, xác định $a^2, b^2 \Rightarrow a, b$ ($a > b > 0$).
\item Tính $c = \sqrt{a^2 - b^2}$[cite: 1].
\item Suy ra các yếu tố: trục lớn $2a$, trục nhỏ $2b$, tiêu cự $2c$, tiêu điểm $F_1(-c; 0), F_2(c; 0)$[cite: 1].
\end{enumerate}
Lưu ý: Nếu phương trình chưa đúng dạng chuẩn (ví dụ $4x^2 + 9y^2 = 36$), phải chia cả 2 vế cho 36 để đưa về dạng chuẩn. $a$ luôn lớn hơn $b$.

%% DANG 0H9?5-2
\textbf{Dạng 2. Phương trình chính tắc của elip}
\begin{enumerate}
\item Gọi phương trình là $\frac{x^2}{a^2} + \frac{y^2}{b^2} = 1$ ($a > b > 0$).
\item Dựa vào giả thiết (tiêu điểm, đi qua điểm, độ dài các trục...) lập hệ $2$ phương trình chứa $a, b, c$.
\item Dùng hệ thức $a^2 = b^2 + c^2$ để giải hệ, tìm $a^2$ và $b^2$[cite: 1].
\end{enumerate}
Lưu ý: Điểm $M(x_0; y_0) \in (E) \Leftrightarrow \frac{x_0^2}{a^2} + \frac{y_0^2}{b^2} = 1$[cite: 1]. Chỉ cần tìm $a^2, b^2$, không cần tính $a, b$ nếu không hỏi.

%% DANG 0H9?5-3
\textbf{Dạng 3. Bài toán điểm trên elip}
\begin{enumerate}
\item Tham số hóa hoặc dùng tọa độ $M(x_0; y_0) \in (E) \Rightarrow \frac{x_0^2}{a^2} + \frac{y_0^2}{b^2} = 1$[cite: 1].
\item Dùng công thức bán kính qua tiêu: $MF_1 = a + \frac{c}{a}x_0; MF_2 = a - \frac{c}{a}x_0$.
\item Kết hợp điều kiện bài toán (góc vuông, khoảng cách, diện tích...) lập phương trình giải tìm $x_0, y_0$.
\end{enumerate}
Lưu ý: Chú ý điều kiện $-a \le x_0 \le a$ và $-b \le y_0 \le b$. Bài toán góc vuông $\widehat{F_1MF_2} = 90^\circ$ thì $MF_1^2 + MF_2^2 = (2c)^2$ hoặc $\overrightarrow{MF_1} \cdot \overrightarrow{MF_2} = 0$.

%% DANG 0H9?5-4
\textbf{Dạng 4. Xác định các yếu tố của hypebol}
\begin{enumerate}
\item Từ phương trình $\frac{x^2}{a^2} - \frac{y^2}{b^2} = 1$, xác định $a^2, b^2 \Rightarrow a, b$ ($a, b > 0$)[cite: 1].
\item Tính $c = \sqrt{a^2 + b^2}$[cite: 1].
\item Suy ra các yếu tố: trục thực $2a$, trục ảo $2b$, tiêu cự $2c$, tiêu điểm $F_1(-c; 0), F_2(c; 0)$[cite: 1].
\end{enumerate}
Lưu ý: Ở hypebol, dấu giữa $x^2$ và $y^2$ là dấu TRỪ. Hypebol có $c^2 = a^2 + b^2$, khác với elip là $c^2 = a^2 - b^2$[cite: 1]. Không có quy định $a$ phải lớn hơn $b$.

%% DANG 0H9?5-5
\textbf{Dạng 5. Phương trình chính tắc của hypebol}
\begin{enumerate}
\item Gọi phương trình là $\frac{x^2}{a^2} - \frac{y^2}{b^2} = 1$ ($a, b > 0$)[cite: 1].
\item Dựa vào giả thiết lập hệ phương trình chứa $a, b, c$.
\item Dùng hệ thức $c^2 = a^2 + b^2$[cite: 1] để giải tìm $a^2, b^2$.
\end{enumerate}
Lưu ý: Hypebol đi qua $M(x_0; y_0)$ thì $\frac{x_0^2}{a^2} - \frac{y_0^2}{b^2} = 1$. Độ dài trục thực là $2a$, trục ảo là $2b$.

%% DANG 0H9?5-6
\textbf{Dạng 6. Bài toán điểm trên hypebol}
\begin{enumerate}
\item Dùng tọa độ $M(x_0; y_0) \in (H) \Rightarrow \frac{x_0^2}{a^2} - \frac{y_0^2}{b^2} = 1$[cite: 1].
\item Dùng định nghĩa: $|MF_1 - MF_2| = 2a$[cite: 1].
\item Lập phương trình theo yêu cầu bài toán để giải tìm $x_0, y_0$.
\end{enumerate}
Lưu ý: Chú ý điều kiện $x_0 \le -a$ (nhánh trái) hoặc $x_0 \ge a$ (nhánh phải).

%% DANG 0H9?5-7
\textbf{Dạng 7. Xác định các yếu tố của parabol}
\begin{enumerate}
\item Từ phương trình $y^2 = 2px$, cho hệ số của $x$ bằng $2p$ để tìm $p$ ($p > 0$)[cite: 1].
\item Tiêu điểm $F\left(\frac{p}{2}; 0\right)$[cite: 1].
\item Phương trình đường chuẩn $\Delta: x = -\frac{p}{2}$[cite: 1].
\end{enumerate}
Lưu ý: Phương trình parabol dạng chuẩn chỉ có $y^2$, hệ số $x$ bậc 1 là số dương. Ví dụ $y^2 = 4x$ thì $2p = 4 \Rightarrow p = 2$.

%% DANG 0H9?5-8
\textbf{Dạng 8. Phương trình chính tắc của parabol}
\begin{enumerate}
\item Gọi phương trình là $y^2 = 2px$ ($p > 0$)[cite: 1].
\item Dựa vào giả thiết (tiêu điểm, đường chuẩn, đi qua điểm) tìm $p$.
\item Viết phương trình.
\end{enumerate}
Lưu ý: Khoảng cách từ tiêu điểm đến đường chuẩn bằng đúng $p$.

%% DANG 0H9?5-9
\textbf{Dạng 9. Bài toán điểm trên parabol}
\begin{enumerate}
\item Gọi $M(x_0; y_0) \in (P) \Rightarrow y_0^2 = 2px_0$.
\item Dùng tính chất $M \in (P) \Leftrightarrow MF = d(M, \Delta) = x_0 + \frac{p}{2}$ (với $\Delta$ là đường chuẩn)[cite: 1].
\item Lập phương trình giải tìm $x_0$.
\end{enumerate}
Lưu ý: Điều kiện $x_0 \ge 0$. Khi có hoành độ $x_0$, thay vào $(P)$ thường tìm được $2$ tung độ $y_0$ đối nhau.

%% DANG 0H9?5-0
\textbf{Dạng 10. Bài toán tổng hợp/thực tế, PP tọa độ hóa 3 đường conic}
\begin{enumerate}
\item Chọn hệ trục $Oxy$ phù hợp, thường chọn gốc $O$ là tâm đối xứng (với Elip/Hypebol) hoặc đỉnh (với Parabol).
\item Xác định tọa độ các điểm biên, điểm quan trọng trên hình vẽ thực tế.
\item Lập phương trình đường conic và dùng phương trình đó để tính toán chiều cao, khoảng cách theo yêu cầu.
\end{enumerate}
Lưu ý: Cần quy đổi đơn vị thực tế (m, km...) tương ứng với các đơn vị trên hệ trục tọa độ $Oxy$.
